
In this thesis, we have taken a comprehensive look at many aspect of p53 dynamics.

We looked at the dynamics of the Y220C mutant and the restorative capabilities of the PK11000 mutation across the p53 splice variants, and found that there was widespread rescue displayed by PK11000 across the different isoforms across global and local dynamics, hydrogen bonding with the bound DNA, and the local energetic network.

Looking at p53 dynamics when bound to different DNA sequences, we saw that both p53 altered DNA dynamics and the DNA sequence altered p53 dynamics. We also hypothesized that there might be a correlation with the strength of p53-DNA interactions and the difference between the DNA sequence and the p53 consensus sequence.

When we looked at the different p53 hotspot mutations, we saw that, through the metrics used to analyze the preliminary data, many mutations do not have drastic differences from the WT p53. We also saw that there were several mutations, such as V157F, that seem to be rescued by PK11000.

Finally, we looked at the restorative compound Rezatapopt when bound to the Y220C p53. We saw that by nearly all metrics tested, restoration was not observed. Most noticeably, the binding of Rezatapopt dramatically altered the C-terminal flexibility and caused a great increase in hydrogen bonding in the C-terminal domain with the DNA.

These studies represent a dramatic increase in our computational understanding of the p53 protein. However, they are limited to our computers. We also set out to design an assay to quantitatively measure p53-depending transcription, which we did by building a system where p53 binds to its binding sequence and promotes the transcription of the \textit{lacZ} gene, which can be easily measured via a colorimetric assay.

The ability to measure p53-dependent transcription immediately makes a variety of experimental questions accessible to us, such as,

\begin{enumerate}
\item Are the different p53 isoforms able to initiate transcription of genes? How does the loss of the transactivation domains, proline-rich domain, or tetramerization effect p53-DNA interactions?
\item How strongly will p53 interact with a wide variety of response elements?
\item Is there a correlation between the similarity of a DNA sequence with the p53 consensus sequence and the strength with which p53 interacts with it?
\item How will mutations alter p53 binding to its response elements?
\item When p53 is mutated, for which sequences will it lose binding affinity for first? Why?
\item Can we treat our yeast cells with small molecules, such as PK11000 or Rezatapopt, and measure a restoration in wild-type-like interactions with the DNA in the Y220C mutant p53, or any other mutation?
\end{enumerate}

It is my sincere hope that in the years to come, these questions will be explored in depth.

To you, who took the time to read this thesis, thank you.